\chapter{Learning on Graphs}
\label{ch:learning_on_graphs}

\section{Introduction to Graph Learning}
Many real-world datasets naturally exhibit a graph or network structure, representing entities and their pairwise interactions. Classical examples include social networks, citation networks, biomedical networks, the Internet, and networks of neurons. 
When working with such data, classical Machine Learning (ML) tasks can be adapted to network structures:
\begin{itemize}
    \item \textbf{Node Classification}: Predicting the type or category of a given node (e.g., classifying the function of proteins in the interactome).
    \item \textbf{Link Prediction}: Predicting whether two nodes are likely to be linked (e.g., content recommendation in social platforms).
    \item \textbf{Community Detection}: Identifying densely linked clusters of nodes.
    \item \textbf{Network Similarity}: Evaluating how similar two (sub)networks are.
\end{itemize}

The main difficulty in applying deep learning to graphs lies in the non-Euclidean geometry of the space. While sounds and texts can be treated as 1D grids, and images as 2D grids, graphs lack a regular spatial structure, making the direct application of standard convolutional filters impossible without a tailored mathematical framework.

\begin{insightbox}[Taxonomy of Graph Learning]
Graph learning techniques can be broadly divided based on the task (Node Classification vs. Graph Classification). For Node Classification, approaches range from \textit{Embedding Techniques} (DeepWalk, LINE, Node2Vec) to \textit{Convolutional Networks} (GCN, GraphSAGE, GAT). For Graph Classification, typical tools involve \textit{Graph Kernels} (Random Walk, Weisfeiler-Lehman) or \textit{Graph Spectrum} methods.
\end{insightbox}

\section{Node Embedding: From Shallow to Deep Encodings}

\subsection{Shallow Encodings}
The goal of node embedding is to learn an efficient, task-independent feature representation for machine learning in networks. An encoder maps each node $u$ to a low-dimensional vector $\mathbf{z}_u \in \mathbb{R}^d$, while a similarity function specifies how relationships in the vector space map to relationships in the original network:
\begin{equation}
    \text{similarity}(u, v) \approx \mathbf{z}_v^\top \mathbf{z}_u
\end{equation}
In a \textit{shallow encoding} approach, the encoder is merely an embedding-lookup matrix $\mathbf{Z} \in \mathbb{R}^{d \times |\mathcal{V}|}$, and nodes are represented via one-hot encoding $\mathbf{v} \in \mathbb{I}^{|\mathcal{V}|}$:
\begin{equation}
    \text{ENC}(v) = \mathbf{Z}\mathbf{v}
\end{equation}
Each node is assigned a unique embedding vector. The key distinction among shallow methods lies in how they define node similarity (e.g., whether they are connected, share neighbors, or possess similar structural roles).

\subsection{Adjacency-based and Multi-hop Similarity}
In adjacency-based similarity, the similarity function approximates the edge weight between $u$ and $v$ in the original network. The loss function we want to minimize is:
\begin{equation}
    \mathcal{L} = \sum_{(u,v) \in \mathcal{V} \times \mathcal{V}} \left\| \mathbf{z}_u^\top \mathbf{z}_v - \mathbf{A}_{u,v} \right\|^2
\end{equation}
where $\mathbf{A}_{u,v}$ is the weighted adjacency matrix. A major drawback of this approach is its $\mathcal{O}(|\mathcal{V}|^2)$ runtime, although it can be reduced to $\mathcal{O}(|\mathcal{E}|)$ by summing only over non-zero edges. Furthermore, it only considers direct, local connections. 

To overcome this, \textit{Multi-hop Similarity} extends the adjacency matrix to consider $k$-hop node neighbors. In practice, a log-transformed, probabilistic adjacency matrix is used:
\begin{equation}
    \tilde{\mathbf{A}}_{i,j}^k = \max \left( \log \left( \frac{(\mathbf{A}_{i,j} / d_i)}{\sum_{l \in \mathcal{V}} (\mathbf{A}_{l,j} / d_l)^k} \right)^k - \alpha, 0 \right)
\end{equation}

\subsection{Random Walk Approaches}
Instead of relying on deterministic similarity metrics, Random Walk approaches interpret the dot product in the feature space as the probability that $u$ and $v$ co-occur on a random walk over the network. 
The embeddings are optimized to maximize the likelihood of random walk co-occurrences. Given a random walk strategy $R$, the loss function is defined as:
\begin{equation}
    \mathcal{L} = \sum_{u \in \mathcal{V}} \sum_{v \in N_R(u)} - \log \left( P(v | \mathbf{z}_u) \right)
\end{equation}
where the probability $P(v | \mathbf{z}_u)$ is typically given by the softmax over all nodes:
\begin{equation}
    P(v | \mathbf{z}_u) = \frac{\exp(\mathbf{z}_u^\top \mathbf{z}_v)}{\sum_{n \in \mathcal{V}} \exp(\mathbf{z}_u^\top \mathbf{z}_n)}
\end{equation}
Computing the nested sum over all nodes results in an $\mathcal{O}(|\mathcal{V}|^2)$ complexity. This is alleviated through \textbf{Negative Sampling}, normalizing against $k$ random "negative" examples instead of the entire graph:
\begin{equation}
    \log \left( \frac{\exp(\mathbf{z}_u^\top \mathbf{z}_v)}{\sum_{n \in \mathcal{V}} \exp(\mathbf{z}_u^\top \mathbf{z}_n)} \right) \approx \log(\sigma(\mathbf{z}_u^\top \mathbf{z}_v)) - \sum_{i=1}^k \log(\sigma(\mathbf{z}_u^\top \mathbf{z}_{n_i}))
\end{equation}
where $n_i \sim P_V$ are sampled proportionally to the nodes' degrees.

\begin{insightbox}[Node2Vec]
While DeepWalk uses fixed-length, unbiased random walks, \textbf{node2vec} introduces a biased random walk that trades off between local and global views of the network. It uses a return parameter $p$ to control the likelihood of returning to the previous node, and an in-out parameter $q$ to control the ratio of outward (DFS-like) to inward (BFS-like) exploration.
\end{insightbox}

\section{Graph Neural Networks (GNNs)}
Shallow encodings suffer from several drawbacks: the number of parameters grows linearly with $|\mathcal{V}|$, they are inherently transductive (cannot generate embeddings for unseen nodes), and they do not incorporate node features. 
GNNs overcome these by generating node embeddings based on local neighborhoods through a process called \textbf{Neighborhood Aggregation} (or Message Passing). Nodes aggregate information from their neighbors using neural networks, enabling the model to be of arbitrary depth. 

The general formulation for the $k$-th layer embedding of node $v$ is:
\begin{equation}
    \mathbf{h}_v^k = \sigma \left( \mathbf{W}_k \sum_{u \in \mathcal{N}(v)} \frac{\mathbf{h}_u^{k-1}}{|\mathcal{N}(v)|} + \mathbf{B}_k \mathbf{h}_v^{k-1} \right)
\end{equation}
where $\mathbf{h}_v^0 = \mathbf{x}_v$ are the initial node features. The model learns the weight matrices $\mathbf{W}_k$ and $\mathbf{B}_k$, sharing them across all nodes. This allows for an \textbf{inductive} capacity: the model can be trained on one graph and generalized to a completely new one, or generate embeddings "on the fly".

\subsection{Graph Convolutional Networks (GCN)}
GCNs are a variation of neighborhood aggregation that introduces more parameter sharing and down-weights high-degree neighbors via symmetric normalization. The update rule is:
\begin{equation}
    \mathbf{h}_v^k = \sigma \left( \mathbf{W}_k \sum_{u \in \mathcal{N}(v) \cup \{v\}} \frac{\mathbf{h}_u^{k-1}}{\sqrt{|\mathcal{N}(u)| |\mathcal{N}(v)|}} \right)
\end{equation}
This can be computed highly efficiently in a vectorized form. Let $\tilde{\mathbf{A}} = \mathbf{A} + \mathbf{I}_N$ be the adjacency matrix with added self-loops, and $\tilde{\mathbf{D}}_{ii} = \sum_j \tilde{\mathbf{A}}_{ij}$ be its corresponding degree matrix. The layer-wise propagation rule is:
\begin{equation}
    \mathbf{H}^{(l+1)} = \sigma \left( \tilde{\mathbf{D}}^{-1/2} \tilde{\mathbf{A}} \tilde{\mathbf{D}}^{-1/2} \mathbf{H}^{(l)} \mathbf{W}^{(l)} \right)
\end{equation}
This is mathematically related to spectral graph convolutions, acting as a first-order approximation of spectral filters.

\subsection{GraphSAGE}
GraphSAGE extends the neighborhood aggregation concept beyond simple weighted averages by allowing any differentiable aggregation function:
\begin{equation}
    \mathbf{h}_v^k = \sigma \left( \left[ \mathbf{A}_k \cdot \text{AGG}(\{\mathbf{h}_u^{k-1}, \forall u \in \mathcal{N}(v)\}) \right] , \mathbf{B}_k \mathbf{h}_v^{k-1} \right)
\end{equation}
Common aggregators include:
\begin{itemize}
    \item \textbf{Mean aggregator}: Simple element-wise mean.
    \item \textbf{Pool aggregator}: Applies a linear projection followed by an element-wise max-pooling operation.
    \item \textbf{LSTM aggregator}: Applies an LSTM to a random permutation of the neighbors' features.
\end{itemize}

\subsection{Gated Graph Neural Networks (GGNN)}
As models go deeper (e.g., >10 layers), standard GNNs suffer from vanishing or exploding gradients and severe overfitting due to the explosion of parameters. \textbf{Gated Graph Neural Networks (GGNN)} mitigate this by borrowing concepts from Recurrent Neural Networks. They apply a Gated Recurrent Unit (GRU) to update the node states iteratively:
\begin{equation}
    \mathbf{m}_v^k = \mathbf{W} \sum_{u \in \mathcal{N}(v)} \mathbf{h}_u^{k-1}
\end{equation}
\begin{equation}
    \mathbf{h}_v^k = \text{GRU}(\mathbf{h}_v^{k-1}, \mathbf{m}_v^k)
\end{equation}
Notice that the aggregation weight matrix $\mathbf{W}$ does not depend on the layer step $k$. This weight sharing across layers acts as a regularizer, enabling the construction of significantly deeper graph models.

\section{Graph Attention Networks (GAT)}
A fundamental limitation of both GCN and GraphSAGE is that they rely on \textit{isotropic} or structurally fixed weights. They assign uniform or degree-based weights to neighbors, failing to distinguish between highly informative and noisy neighbors based on their actual features. 
GAT solves this by introducing a dynamically learned \textbf{attention mechanism} to compute adaptive, data-dependent weights.

\subsection{The Attention Mechanism}
In GAT, the computation of the new representation $\mathbf{h}_v^k$ involves the following steps:
1. \textbf{Linear Transformation}: Apply a shared learnable weight matrix $\mathbf{W}$ to project node features into a higher-dimensional space: $\mathbf{W}\mathbf{h}_u^{k-1}$.
2. \textbf{Compute Attention Scores}: For each edge $(u,v)$, compute an unnormalized score $e_{vu}$ using a learnable parameter vector $\mathbf{w}_{attn}$:
\begin{equation}
    e_{vu} = \text{LeakyReLU} \left( \mathbf{w}_{attn}^\top [ \mathbf{W}\mathbf{h}_v^{k-1} \| \mathbf{W}\mathbf{h}_u^{k-1} ] \right)
\end{equation}
3. \textbf{Normalization}: Apply softmax over the connected neighborhood $\mathcal{N}(v)$ to obtain the normalized attention coefficient $\alpha_{vu}$:
\begin{equation}
    \alpha_{vu} = \text{softmax}_u(e_{vu}) = \frac{\exp(e_{vu})}{\sum_{k \in \mathcal{N}(v)} \exp(e_{vk})}
\end{equation}
4. \textbf{Weighted Aggregation}: Compute the final feature representation as a weighted sum followed by a non-linearity $\sigma$:
\begin{equation}
    \mathbf{h}_v^k = \sigma \left( \sum_{u \in \mathcal{N}(v)} \alpha_{vu} \mathbf{W}_k \mathbf{h}_u^{k-1} \right)
\end{equation}

\subsection{Multi-Head Attention}
To stabilize learning and capture different facets of importance, GAT uses $I$ independent attention heads computed in parallel. For intermediate layers, the outputs are concatenated:
\begin{equation}
    \mathbf{h}_v^k = \Big\|_{i=1}^I \mathbf{h}_v^{k(i)}
\end{equation}
For the final layer, to ensure the representation matches the desired output dimensions, the outputs are averaged instead:
\begin{equation}
    \mathbf{h}_v^{k(i)} = \sigma \left( \frac{1}{I} \sum_{i=1}^I \sum_{u \in \mathcal{N}(v)} \alpha_{vu}^{(i)} \mathbf{W}_k^{(i)} \mathbf{h}_u^{k-1(i)} \right)
\end{equation}

\begin{insightbox}[Limitation: Static Attention in GAT]
While GAT computes data-dependent weights, the mechanism is \textit{static}: the MLP learns a single fixed ranking function over the projected features. It cannot dynamically adapt the importance of different feature dimensions for different query-key pairs. \textbf{GATv2} addresses this by applying the weight matrix $\mathbf{W}$ \textit{after} the concatenation: $e_{vu} = \mathbf{w}_{attn}^\top \text{LeakyReLU} \left( \mathbf{W} [\mathbf{h}_v^{k-1} \| \mathbf{h}_u^{k-1}] \right)$, achieving universal expressiveness.
\end{insightbox}

\section{Beyond Message Passing: Graph Transformers}
Despite their successes, all Message Passing Neural Networks (MPNNs), including GATv2, share fundamental limitations:
\begin{itemize}
    \item \textbf{Over-smoothing}: Because information propagates 1-hop per layer, capturing $K$-hop relations requires $K$ layers. With large $K$, node representations converge and become indistinguishable.
    \item \textbf{Over-squashing}: A node's entire $K$-hop neighborhood must be compressed ("squashed") into a single fixed-size vector, causing inevitable information loss.
\end{itemize}

\subsection{The Graphormer Approach}
To overcome the locality limits of MPNNs, the \textbf{Graphormer} abandons the strict neighborhood-aggregation paradigm and treats the graph as a set of $N$ nodes, applying the Self-Attention mechanism from Transformers to \textit{all pairs of nodes}:
\begin{equation}
    \text{Attention}(\mathbf{H}) = \text{softmax} \left( \frac{\mathbf{Q}\mathbf{K}^\top}{\sqrt{d_k}} \right) \mathbf{V}
\end{equation}
However, a standard Transformer is permutation-invariant and loses all topological information. Graphormer re-injects the graph structure through several specialized encodings:
\begin{enumerate}
    \item \textbf{Centrality Encoding}: Adds a learned embedding based on a node's degree, capturing its importance.
    \item \textbf{Spatial Encoding}: Adds a learned attention bias $b_{\phi(u,v)}$ based on the Shortest Path Distance (SPD) between node $u$ and $v$.
    \item \textbf{Edge Encoding}: Adds a learned bias $w_{e_{vu}}$ based on the average relationship of edges along the shortest path between $u$ and $v$.
\end{enumerate}
The resulting modified attention score matrix $\mathbf{A}_{vu}$ becomes:
\begin{equation}
    A_{vu} = \frac{(\mathbf{h}_v \mathbf{W}_Q)(\mathbf{h}_u \mathbf{W}_K)^\top}{\sqrt{d}} + b_{\phi(u,v)} + w_{e_{vu}}
\end{equation}
This effectively "teaches" the Transformer about connectivity, distance, and the nature of relations, resolving the MPNN limitations by allowing long-range, direct message passing without over-squashing.

An alternative method for Positional Encoding utilizes the \textbf{Laplacian Eigenvectors} (LapPE), which capture both geometric and spectral structures of the graph by taking the first $k$ eigenvectors of the normalized Graph Laplacian $\mathbf{L}_{frw} = \mathbf{I} - \mathbf{D}^{-1}\mathbf{A}$.

\section{Applications: Spatial-Temporal Graphs and Beyond}

\subsection{Recommender Systems and Polypharmacy}
GCNs and their variants have widespread applications. In \textbf{Recommender Systems}, the data can be structured as a bipartite graph of users and items. The \textit{PinSage} algorithm (used by Pinterest) employs Random Walk Graph Convolutional Networks (RW-GCN) to dynamically generate embeddings for pins "on the fly" without needing full model retraining. It leverages sub-sampled neighborhoods based on top-K PageRank scores for efficient GPU batching and uses a max-margin ranking loss with curriculum learning for negative samples:
\begin{equation}
    \mathcal{L} = \sum_{(u,v) \in \mathcal{D}} \max(0, - \mathbf{z}_u^\top \mathbf{z}_v + \mathbf{z}_u^\top \mathbf{z}_n + \Delta)
\end{equation}
where $(u,v)$ is a positive training pair (consecutively clicked pins), $n$ is a negative sample, and $\Delta$ is the margin.

Another powerful application is predicting \textbf{Polypharmacy Side Effects}. Patients often take multiple drugs simultaneously, leading to rare adverse side effects that were not observed during isolated clinical testing. By framing this as a link prediction problem within an encoder-decoder architecture, nodes representing drugs and proteins are embedded using a GCN, and labeled edges predicting specific side effects $r$ are inferred using a multirelational link prediction decoder:
\begin{equation}
    p(r | u, v) = \sigma( \mathbf{z}_u^\top \mathbf{R}_r \mathbf{z}_v )
\end{equation}
This computational screening guides the formulation of safer combination treatments.

\subsection{Skeleton-based Activity Recognition}
Many real-world graphs are dynamic, evolving over time. These spatial-temporal graphs appear in recommender systems, financial transactions, and notably in video analysis for \textbf{Skeleton-based Activity Recognition}. 
In this domain, human body joints form the nodes of a graph, while bones act as edges. Standard models like the Spatial-Temporal Graph Convolutional Network (ST-GCN) apply spatial convolutions within frames and temporal convolutions across frames. 
However, ST-GCNs suffer from a fixed topology that struggles to capture long-range correlations. A modern solution is the \textbf{Spatial Temporal Transformer Network} (ST-TR), which leverages a Spatial Self-Attention module within each frame, and a Temporal Self-Attention module between frames, allowing for dynamic, long-range correlation modeling both spatially and temporally.

\section{Chapter Overview}

This chapter introduces the fundamental concepts of learning on graphs, transitioning from shallow node embeddings to advanced deep learning architectures. It begins with shallow encodings, such as adjacency-based similarity and random walk approaches (e.g., DeepWalk and node2vec), which are transductive and scale linearly with the number of nodes. To address these limitations, Graph Neural Networks (GNNs) are introduced, utilizing Neighborhood Aggregation (Message Passing) to enable inductive learning and incorporate node features. Key GNN variants discussed include Graph Convolutional Networks (GCNs), GraphSAGE, and Gated Graph Neural Networks (GGNNs). To overcome the static weighting of neighbors in standard GNNs, Graph Attention Networks (GAT) introduce a dynamic, data-dependent attention mechanism. Furthermore, the chapter addresses the inherent limitations of Message Passing Neural Networks (MPNNs)—namely over-smoothing and over-squashing—by exploring Graph Transformers (like Graphormer), which apply self-attention globally and re-inject graph topology via specialized positional and structural encodings. Finally, the chapter highlights practical applications, including recommender systems (PinSage), polypharmacy, and skeleton-based activity recognition.

\begin{table}[h]
\centering
\begin{tabularx}{\textwidth}{| >{\centering\arraybackslash}m{3.5cm} | X | X |}
\hline
\textbf{Architecture / Technique} & \textbf{Pros} & \textbf{Cons} \\
\hline
\textbf{Shallow Encodings (e.g., node2vec)} & Efficient for local similarities; well-suited for transductive tasks & Parameters grow as $\mathcal{O}(|\mathcal{V}|)$; purely transductive; ignores node features \\
\hline
\textbf{Graph Convolutional Networks (GCN)} & Incorporates node features; inductive capability; efficient vectorized computation & Relies on isotropic/static weights; suffers from over-smoothing in deep models \\
\hline
\textbf{GraphSAGE} & Flexible aggregation (Mean, Pool, LSTM); effective for unseen nodes & Inherits locality constraints of standard MPNNs \\
\hline
\textbf{Gated Graph Neural Networks (GGNN)} & Alleviates overfitting and vanishing gradients, allowing deeper models (>10 layers) & Recurrent GRU updates increase computational overhead \\
\hline
\textbf{Graph Attention Networks (GAT)} & Dynamic, data-dependent neighbor weighting; scalable with multi-head attention & Standard GAT attention is static/limited expressiveness (solved by GATv2) \\
\hline
\textbf{Graph Transformers (Graphormer)} & Captures long-range dependencies; effectively mitigates over-smoothing and over-squashing & $\mathcal{O}(|\mathcal{V}|^2)$ complexity; requires complex structural/spatial encodings \\
\hline
\end{tabularx}
\end{table}

\subsection{Possible Exam Questions}
\begin{itemize}
    \item \textbf{Q: What are the main limitations of shallow encoding approaches for node embedding?} \\
    \textbf{A:} They are transductive (cannot generate embeddings for unseen nodes), the number of parameters grows linearly with the number of nodes $\mathcal{O}(|\mathcal{V}|)$, and they ignore node features.
    \item \textbf{Q: Explain the concept of Neighborhood Aggregation (Message Passing) and how it enables inductive learning in GNNs.} \\
    \textbf{A:} Nodes aggregate information from their local neighbors using neural networks and shared weight matrices. Because the weights are shared across all nodes rather than tied to specific nodes, the model can generate embeddings for unseen nodes or graphs on the fly.
    \item \textbf{Q: Describe how random walk approaches like DeepWalk or node2vec optimize node embeddings through Negative Sampling.} \\
    \textbf{A:} Negative Sampling reduces the computational cost of the softmax loss from $\mathcal{O}(|\mathcal{V}|^2)$ to $\mathcal{O}(k)$ by normalizing against $k$ randomly sampled "negative" nodes (drawn proportionally to their degree) instead of the entire graph.
    \item \textbf{Q: Why do standard message-passing models like GCNs struggle when scaled to many layers, and what is this phenomenon called?} \\
    \textbf{A:} It is called over-smoothing. Since information propagates 1-hop per layer, adding many layers causes node representations to repeatedly mix, eventually converging to similar values and becoming indistinguishable.
    \item \textbf{Q: How does the aggregation mechanism in Graph Attention Networks (GAT) differ fundamentally from that in Graph Convolutional Networks (GCN)?} \\
    \textbf{A:} GCN uses isotropic or structurally fixed weights (like degree-based normalization) for neighbors, while GAT introduces a dynamically learned attention mechanism that computes adaptive, data-dependent weights for each neighbor.
    \item \textbf{Q: Discuss the "over-squashing" and "over-smoothing" problems in Message Passing Neural Networks (MPNNs).} \\
    \textbf{A:} Over-smoothing occurs when many layers cause node representations to converge and become indistinguishable. Over-squashing happens when a node's entire exponentially growing $K$-hop neighborhood must be compressed into a single fixed-size vector, causing information loss.
    \item \textbf{Q: How does the Graphormer inject graph topological information into the permutation-invariant standard Transformer architecture?} \\
    \textbf{A:} It uses specialized encodings: Centrality Encoding (node degree), Spatial Encoding (Shortest Path Distance between nodes), and Edge Encoding (average relationship of edges along the shortest path).
    \item \textbf{Q: Explain the difference between transductive and inductive learning in the context of graph networks.} \\
    \textbf{A:} Transductive learning requires the entire graph structure (including test nodes) to be present during training and cannot generate embeddings for new, unseen nodes. Inductive learning learns mapping functions (e.g., via shared weights) that can generalize to entirely new nodes or graphs.
\end{itemize}
